% ============================================================================
% CONCLUSIONES
% ============================================================================
% Las conclusiones constituyen la síntesis de los logros obtenidos en el
% trabajo. Deben responder a los objetivos planteados en la introducción.
%
% Según la norma UTEM, deben abarcar:
%   a) Evaluación del cumplimiento de los objetivos (general y específicos).
%   b) Aportaciones del trabajo al área de conocimiento.
%   c) Limitaciones encontradas.
%   d) Trabajos futuros (propuestas de desarrollo o investigación posterior).
% ============================================================================

\chapter*{CONCLUSIONES}
\addcontentsline{toc}{chapter}{CONCLUSIONES}

% --- Introducción a las conclusiones ---
% Presenta un breve balance general sobre la ejecución del proyecto...

\section*{Cumplimiento de Objetivos}

[Evalúa críticamente cómo se alcanzaron los objetivos propuestos (tanto el general como los específicos) y qué resultados concretos respaldan este cumplimiento.]

\section*{Aportes y Contribuciones}

[Describe el valor agregado de tu trabajo. ¿Qué aporta a la disciplina, a la institución, a la empresa (si aplica) o al estado del arte?]

\section*{Limitaciones de la Solución}

[Describe con honestidad técnica las restricciones y limitaciones de tu desarrollo o de tu investigación (por ejemplo, dependencias tecnológicas, falta de datos, limitaciones de rendimiento, etc.).]

\section*{Trabajos Futuros}

[Propón sugerencias, recomendaciones o nuevas líneas de trabajo, mejoras de software y extensiones de la investigación que quedaron fuera del alcance de este proyecto pero que son viables para continuarse en el futuro.]
